% Options for packages loaded elsewhere
\PassOptionsToPackage{unicode}{hyperref}
\PassOptionsToPackage{hyphens}{url}
\PassOptionsToPackage{dvipsnames,svgnames,x11names}{xcolor}
\documentclass[
  11pt,
]{article}
\usepackage{xcolor}
\usepackage[a4paper,margin=19mm]{geometry}
\usepackage{amsmath,amssymb}
\setcounter{secnumdepth}{-\maxdimen} % remove section numbering
\usepackage{iftex}
\ifPDFTeX
  \usepackage[T1]{fontenc}
  \usepackage[utf8]{inputenc}
  \usepackage{textcomp} % provide euro and other symbols
\else % if luatex or xetex
  \usepackage{unicode-math} % this also loads fontspec
  \defaultfontfeatures{Scale=MatchLowercase}
  \defaultfontfeatures[\rmfamily]{Ligatures=TeX,Scale=1}
\fi
\usepackage{lmodern}
\ifPDFTeX\else
  % xetex/luatex font selection
  \setmainfont[]{Be Vietnam Pro}
  \setmonofont[]{Menlo}
\fi
% Use upquote if available, for straight quotes in verbatim environments
\IfFileExists{upquote.sty}{\usepackage{upquote}}{}
\IfFileExists{microtype.sty}{% use microtype if available
  \usepackage[]{microtype}
  \UseMicrotypeSet[protrusion]{basicmath} % disable protrusion for tt fonts
}{}
\makeatletter
\@ifundefined{KOMAClassName}{% if non-KOMA class
  \IfFileExists{parskip.sty}{%
    \usepackage{parskip}
  }{% else
    \setlength{\parindent}{0pt}
    \setlength{\parskip}{6pt plus 2pt minus 1pt}}
}{% if KOMA class
  \KOMAoptions{parskip=half}}
\makeatother
\usepackage{longtable,booktabs,array}
\usepackage{caption}
\captionsetup[table]{skip=6pt}
\newcounter{none} % for unnumbered tables
\usepackage{calc} % for calculating minipage widths
% Correct order of tables after \paragraph or \subparagraph
\usepackage{etoolbox}
\makeatletter
\patchcmd\longtable{\par}{\if@noskipsec\mbox{}\fi\par}{}{}
\makeatother
% Allow footnotes in longtable head/foot
\IfFileExists{footnotehyper.sty}{\usepackage{footnotehyper}}{\usepackage{footnote}}
\makesavenoteenv{longtable}
\setlength{\emergencystretch}{3em} % prevent overfull lines
\providecommand{\tightlist}{%
  \setlength{\itemsep}{0pt}\setlength{\parskip}{0pt}}

\usepackage{newunicodechar}
\newunicodechar{→}{\ensuremath{\rightarrow}}
\newunicodechar{↔}{\ensuremath{\leftrightarrow}}
\usepackage{fvextra}
\DefineVerbatimEnvironment{verbatim}{Verbatim}{breaklines=true,breakanywhere=true,fontsize=\small}
\usepackage{fancyhdr}
\usepackage{array}
\pagestyle{fancy}\fancyhf{}
\fancyhead[L]{\small MovieLOD · Bài làm hoàn chỉnh}
\fancyfoot[C]{\small\thepage}\setlength{\headheight}{15pt}
\renewcommand{\arraystretch}{1.14}\setlength{\extrarowheight}{1pt}
\AtBeginEnvironment{longtable}{\small}
\setlength{\emergencystretch}{3em}\setlength{\parskip}{4pt}
\usepackage{titling}\setlength{\droptitle}{-1.5cm}
\pretitle{\begin{center}\Large\bfseries\color{teal}}
\posttitle{\par\end{center}\vspace{-0.5em}}
\predate{\begin{center}\small}\postdate{\par\end{center}\vspace{-1em}}
\usepackage{bookmark}
\IfFileExists{xurl.sty}{\usepackage{xurl}}{} % add URL line breaks if available
\urlstyle{same}
% fallback for those not using the hyperref driver hyperxmp:
\makeatletter
\@ifundefined{xmpquote}{\newcommand{\xmpquote}[1]{#1}}{}
\makeatother
\hypersetup{
  pdftitle={MovieLOD: lời thuyết trình cho bộ 24 slide},
  colorlinks=true,
  linkcolor={teal},
  filecolor={Maroon},
  citecolor={Blue},
  urlcolor={teal},
  pdfcreator={LaTeX via pandoc}}

\title{MovieLOD: lời thuyết trình cho bộ 24 slide}
\author{}
\date{Bản đầy đủ · 08/10/2026}

\begin{document}
\maketitle

\subsection{Cách dùng}\label{cuxe1ch-duxf9ng}

Bộ chính có \textbf{24 slide}, thuyết trình đầy đủ khoảng \textbf{18--22
phút}. Nội dung slide và Speaker Notes bằng tiếng Anh; script riêng giữ
lời tiếng Việt tương ứng để tập nói. Hướng dẫn trong ô chờ ảnh vẫn bằng
tiếng Việt. Bản ngắn 13 trang vẫn ở
Slide\_\allowbreak{}ng\allowbreak{}an\_\allowbreak{}13.\allowbreak{}pp\allowbreak{}tx/\allowbreak{}pdf
và
Script\_\allowbreak{}t\allowbreak{}huyet\_\allowbreak{}tr\allowbreak{}inh\_\allowbreak{}ngan\allowbreak{}\_\allowbreak{}13.\allowbreak{}md/\allowbreak{}p\allowbreak{}df.\allowbreak{}
Video demo 4:50 là tài liệu riêng; không đọc toàn bộ lời slide vào
video.

Tập theo 3 phần: thành viên A slide 1--10; B slide 11--18; C slide
19--24. Nếu chỉ có 2 người, chia sau slide 14. Các con số lấy từ dữ liệu
và evidence hiện tại. Khung ảnh Protégé chưa có ảnh thực; không đọc ghi
chú chờ bổ sung như kết quả đã chứng minh.

\subsection{Luồng rút gọn 12--15
phút}\label{luux1ed3ng-ruxfat-gux1ecdn-1215-phuxfat}

Ưu tiên slide 1--5, 9--10, 12--13, 15--17, 19--23. Các trang cây lớp,
công thức cardinality và tra cứu IRI có thể dùng khi trả lời câu hỏi.

\subsection{Mục lục}\label{mux1ee5c-lux1ee5c}

{\def\LTcaptype{none} % do not increment counter
\begin{longtable}[]{@{}ll@{}}
\toprule\noalign{}
Slide & Nội dung \\
\midrule\noalign{}
\endhead
\bottomrule\noalign{}
\endlastfoot
01 & MovieLOD \\
02 & Problem and five assignment requirements \\
03 & Architecture and data flow \\
04 & Current dataset: interpreting the numbers \\
05 & R1 · Ontology: 42 classes in six groups \\
06 & R1 · Creative works and film classes \\
07 & R1 · People, organiza\allowbreak{}tions and roles \\
08 & R1 · Genres, awards and provenance \\
09 & R1 · Contribu\allowbreak{}tion:\allowbreak{} who did what in a
film? \\
10 & R1 · Properties, domain/range and inverses \\
11 & R1 · How OWL expresses meaning \\
12 & R1 · Fourteen classes with additional types \\
13 & R1 · Inference example: Christopher Nolan \\
14 & R1 · Cardinality: thresholds of two and three \\
15 & R2 · Real collection and traceable sources \\
16 & R3 · RDF and typed literal values \\
17 & R3--R4 · From RDF to Linked Open Data \\
18 & R3 · Looking up Inception's IRI \\
19 & R5 · Querying Inception on the Web \\
20 & R5 · One person, multiple roles and questions \\
21 & R5 · Endpoint, terminal and inferred types \\
22 & Validating data, application and publication \\
23 & Assessment and limitations \\
24 & Conclusion and pending Protégé evidence \\
\end{longtable}
}

\subsection{Slide 01 --- MovieLOD}\label{slide-01-movielod}

\textbf{Lời nói:}

Chào thầy cô và các bạn. Nhóm chúng em trình bày MovieLOD, ứng dụng dữ
liệu mở có liên kết về điện ảnh. Bài đi từ mô hình khái niệm, dữ liệu
nguồn và RDF đến liên kết ngoài và truy vấn. Phiên bản hiện tại có 30
phim, 42 lớp ontology và 1.727 liên kết ngoài. Ví dụ xuyên suốt là
Inception và Christopher Nolan. Bộ slide này có 24 trang; lời nói riêng
và hướng dẫn chụp minh chứng Protégé đi kèm. Một số khung ảnh Protégé
được để sẵn vì lần làm tài liệu này chưa chụp được giao diện trên máy.
Các khung này ghi rõ yêu cầu, không được coi là ảnh minh chứng đã có.

\textbf{Chuyển trang:} Sau đây nhóm chuyển sang problem and five
assignment requirem\allowbreak{}ents.\allowbreak{}

\subsection{Slide 02 --- Problem and five assignment
requirements}\label{slide-02-problem-and-five-assignment-requirements}

\textbf{Lời nói:}

Điện ảnh có nhiều thực thể và nhiều vai trò chồng lấp. Chúng ta muốn hỏi
Nolan tham gia Inception ở những vai trò nào, phim này liên quan tới
công ty nào hoặc dữ liệu lấy từ phản hồi nào. Năm yêu cầu của đề tạo
thành một chuỗi công việc. Ontology định nghĩa khái niệm và quan hệ. Thu
thập cung cấp dữ liệu thật. Chuyển đổi tạo RDF và IRI với giấy phép mở.
Liên kết nối cá thể tới Wikidata và DBpedia. SPARQL cho phép đặt câu hỏi
trên đồ thị. Đề còn yêu cầu báo cáo không quá 15 trang và video 3--5
phút; video demo hiện dài khoảng 4 phút 50 giây.

\textbf{Chuyển trang:} Sau đây nhóm chuyển sang architecture and data
flow.

\subsection{Slide 03 --- Architecture and data
flow}\label{slide-03-architecture-and-data-flow}

\textbf{Lời nói:}

Luồng xử lý bắt đầu từ danh sách phim trong config. collect lấy phản hồi
nguồn và lưu metadata. build tạo ontology cùng graph RDF từ dữ liệu đã
chuẩn hóa. validate kiểm tra trường ứng dụng, hash và truy vấn. reason
bổ sung các kiểu phân loại ở file riêng. Web cục bộ gọi endpoint Flask
dùng RDFLib; hosted dùng Comunica trong trình duyệt. Endpoint mặc định
đọc movies.ttl, còn terminal với reasoned nạp thêm schema và
inferred\allowbreak{}\_\allowbreak{}classes\allowbreak{}.\allowbreak{}ttl.\allowbreak{}
Tùy chọn này nạp kết quả phân loại đã lưu, không tự chạy reasoner mỗi
lần truy vấn. Việc tách ba phần giúp phân biệt dữ kiện gốc với kiểu được
bổ sung.

\textbf{Chuyển trang:} Sau đây nhóm chuyển sang current dataset:
interpreting the numbers.

\subsection{Slide 04 --- Current dataset: interpreting the
numbers}\label{slide-04-current-dataset-interpreting-the-numbers}

\textbf{Lời nói:}

Mẫu hiện có 30 phim, 851 người và 1.010 bản ghi đóng góp. Ngoài ra có 45
công ty, 672 thực thể giải thưởng, 75 thể loại, 11 quốc gia và 15 ngôn
ngữ. 76 phản hồi gốc được giữ và kiểm tra hash. Dữ liệu gồm 19.339
triple, tách khỏi 518 triple lược đồ. Cần đọc đúng đơn vị: 672 là số
thực thể giải khác nhau, không phải số lần trao giải; 1.010 là bản ghi
một người giữ một vai trò trong một phim, không phải số người. Mọi phim
trong mẫu có năm, thời lượng và đạo diễn, nhưng mẫu này không đại diện
toàn bộ điện ảnh.

\textbf{Chuyển trang:} Sau đây nhóm chuyển sang r1 · ontology: 42
classes in six groups.

\subsection{Slide 05 --- R1 · Ontology: 42 classes in six
groups}\label{slide-05-r1-ontology-42-classes-in-six-groups}

\textbf{Lời nói:}

Ontology có 42 IRI lớp được khai báo owl Class, trong đó ba lớp dùng
trực tiếp từ DBpedia, còn lại thuộc namespace của bài. Chúng ta có nhánh
tác phẩm và phim; chủ thể người hoặc tổ chức; đóng góp và vai trò; thể
loại; giải thưởng; và các lớp bổ trợ cho quốc gia, ngôn ngữ, nguồn,
dataset. Trong số này, 14 lớp có kiểu được bổ sung từ định nghĩa hoặc
quy tắc phân loại. Các lớp ngoài được tham chiếu qua
equivale\allowbreak{}nt\allowbreak{}Class có thể được Protégé tính thêm
trong Metrics; vì vậy con số 42 dùng quy ước đếm lớp tự khai báo. Ảnh
cần bổ sung ở bên phải là IRI ontology và versionInfo 2.0.0.

\textbf{Ảnh cần bổ sung --- P00:}
\texttt{P00\_\allowbreak{}onto\allowbreak{}logy\_\allowbreak{}hea\allowbreak{}der.\allowbreak{}png};
mở
\texttt{ontology\allowbreak{}/\allowbreak{}\allowbreak{}Movie\_\allowbreak{}\allowbreak{}O\allowbreak{}ntology.\allowbreak{}owl}.

\begin{enumerate}
\def\labelenumi{\arabic{enumi}.}
\tightlist
\item
  Mở
  ontology\allowbreak{}/\allowbreak{}\allowbreak{}Movie\_\allowbreak{}\allowbreak{}O\allowbreak{}ntology.\allowbreak{}owl.\allowbreak{}
\item
  Chọn Active Ontology / Ontology Header.
\item
  Giữ IRI ontology và versionInfo 2.0.0 trong khung hình.
\end{enumerate}

\textbf{Cần thấy:} IRI đúng namespace của bài; phiên bản 2.0.0. Metrics
có thể tính cả lớp ngoài được tham chiếu, khác 42 lớp tự khai báo.

\textbf{Chuyển trang:} Sau đây nhóm chuyển sang r1 · creative works and
film classes.

\subsection{Slide 06 --- R1 · Creative works and film
classes}\label{slide-06-r1-creative-works-and-film-classes}

\textbf{Lời nói:}

CreativeWork là nhánh tác phẩm. Film dùng IRI DBpedia và có các lớp con
như FeatureFilm, Animated\allowbreak{}\allowbreak{}Film,
Document\allowbreak{}ary\allowbreak{}Film.\allowbreak{} FeatureFilm và
AnimatedFilm được ánh xạ theo quy tắc P31 của bài; đây là cách giản
lược, không phải khẳng định đã kiểm tra mọi tiêu chí phát hành rạp. Các
lớp ActionFilm, ComedyFilm, DramaFilm,
Science\allowbreak{}F\allowbreak{}iction\allowbreak{}Fi\allowbreak{}lm,
Multi\allowbreak{}Gen\allowbreak{}re\allowbreak{}Film và
Award\allowbreak{}Win\allowbreak{}ning\allowbreak{}Film\allowbreak{}
được phân loại bổ sung. Document\allowbreak{}ary\allowbreak{}Film hiện
không có cá thể trong mẫu. Việc có một lớp trong mô hình không yêu cầu
mẫu dữ liệu phải có sẵn mọi loại phim. Khung ảnh yêu cầu mở cây
Creative\allowbreak{}\allowbreak{}Work/\allowbreak{}\allowbreak{}Fil\allowbreak{}m
và chọn Film để thấy IRI DBpedia.

\textbf{Ảnh cần bổ sung --- P01:}
\texttt{P01\_\allowbreak{}film\allowbreak{}\_\allowbreak{}hierarc\allowbreak{}hy.\allowbreak{}png};
mở
\texttt{ontology\allowbreak{}/\allowbreak{}\allowbreak{}Movie\_\allowbreak{}\allowbreak{}O\allowbreak{}ntology.\allowbreak{}owl}.

\begin{enumerate}
\def\labelenumi{\arabic{enumi}.}
\tightlist
\item
  Entities → Classes → mở CreativeWork và Film.
\item
  Hiện các lớp con FeatureFilm, AnimatedFilm và nhóm phim.
\item
  Chọn Film; giữ IRI DBpedia và Class Description.
\end{enumerate}

\textbf{Cần thấy:} Film thuộc Creative\allowbreak{}\allowbreak{}Work;
IRI là
http:\allowbreak{}/\allowbreak{}/\allowbreak{}d\allowbreak{}bpedia.\allowbreak{}o\allowbreak{}rg/\allowbreak{}ontol\allowbreak{}ogy/\allowbreak{}\allowbreak{}Film\allowbreak{}.\allowbreak{}
Chụp cây khai báo, không gọi là cây đã suy luận.

\textbf{Chuyển trang:} Sau đây nhóm chuyển sang r1 · people,
organiza\allowbreak{}tions and roles.

\subsection{\texorpdfstring{Slide 07 --- R1 · People,
organiza\allowbreak{}tions and
roles}{Slide 07 --- R1 · People,  and roles}}\label{slide-07-r1-people-organizations-and-roles}

\textbf{Lời nói:}

Agent gồm Person và Organiza\allowbreak{}tion.\allowbreak{} Person dùng
lớp DBpedia và có các nhóm Actor, Filmmaker, AwardWinner. Organization
có Producti\allowbreak{}on\allowbreak{}Compan\allowbreak{}y; FilmStudio
là nhóm công ty được ứng dụng phân loại theo ngưỡng số phim trong mẫu.
Actor và Filmmaker không rời nhau, vì một người có thể vừa diễn xuất vừa
đạo diễn hoặc biên kịch. Chúng ta cũng không dùng một lớp tổ hợp
Director\allowbreak{}\allowbreak{}Writer để giữ hai nghề; từng công việc
được ghi qua Contribution và SPARQL có thể JOIN các đóng góp. Số người
hoặc công ty không được suy từ số dòng của một truy vấn có LIMIT. Ảnh
cần chụp là cây Agent mở cả nhánh Person và
Organiza\allowbreak{}tion.\allowbreak{}

\textbf{Ảnh cần bổ sung --- P02:}
\texttt{P02\_\allowbreak{}agen\allowbreak{}t\_\allowbreak{}hierar\allowbreak{}chy.\allowbreak{}png};
mở
\texttt{ontology\allowbreak{}/\allowbreak{}\allowbreak{}Movie\_\allowbreak{}\allowbreak{}O\allowbreak{}ntology.\allowbreak{}owl}.

\begin{enumerate}
\def\labelenumi{\arabic{enumi}.}
\tightlist
\item
  Entities → Classes → mở Agent.
\item
  Mở Person và Organization →
  Producti\allowbreak{}on\allowbreak{}Compan\allowbreak{}y.\allowbreak{}
\item
  Chọn Person hoặc Filmmaker; giữ cây lớp và Description.
\end{enumerate}

\textbf{Cần thấy:} Person và Organization cùng dưới Agent;
Actor/\allowbreak{}\allowbreak{}Fi\allowbreak{}lmmaker/\allowbreak{}\allowbreak{}Award\allowbreak{}Win\allowbreak{}ner
là lớp con Person.

\textbf{Chuyển trang:} Sau đây nhóm chuyển sang r1 · genres, awards and
provenance.

\subsection{Slide 08 --- R1 · Genres, awards and
provenance}\label{slide-08-r1-genres-awards-and-provenance}

\textbf{Lời nói:}

Genre có FictionGenre và
Non\allowbreak{}Ficti\allowbreak{}on\allowbreak{}Genre;\allowbreak{} các
lớp thể loại hành động, hài, chính kịch, khoa học viễn tưởng thuộc
Fiction\allowbreak{}G\allowbreak{}enre, còn
Document\allowbreak{}ary\allowbreak{}Genre\allowbreak{} dưới
Non\allowbreak{}Ficti\allowbreak{}on\allowbreak{}Genre.\allowbreak{}
Award có các nhóm FilmAward, ActingAward,
Directin\allowbreak{}g\allowbreak{}Award,
Writing\allowbreak{}A\allowbreak{}ward.\allowbreak{} Việc chia nhóm thể
loại và giải hiện dùng từ khóa nhãn trong bước build, không phải suy
luận tự khám phá nghĩa tiếng Anh. FilmAward là nhóm mặc định khi nhãn
chưa khớp một nhóm nghề nghiệp. Các lớp Country, Language,
Source\allowbreak{}Sn\allowbreak{}apshot và Dataset bổ sung bối cảnh và
xuất xứ. Giải của người được giữ riêng với giải của phim. Khung ảnh cần
mở hai nhánh Genre và Award, thu gọn nhánh khác để đọc rõ.

\textbf{Ảnh cần bổ sung --- P03:}
\texttt{P03\_\allowbreak{}genr\allowbreak{}e\_\allowbreak{}award\_\allowbreak{}hierarch\allowbreak{}y.\allowbreak{}png};
mở
\texttt{ontology\allowbreak{}/\allowbreak{}\allowbreak{}Movie\_\allowbreak{}\allowbreak{}O\allowbreak{}ntology.\allowbreak{}owl}.

\begin{enumerate}
\def\labelenumi{\arabic{enumi}.}
\tightlist
\item
  Thu gọn Film/Agent để ảnh dễ đọc.
\item
  Mở Genre → FictionGenre /
  Non\allowbreak{}Ficti\allowbreak{}on\allowbreak{}Genre và Award.
\item
  Hiện các lớp thể loại và nhóm giải; giữ Class Description.
\end{enumerate}

\textbf{Cần thấy:} Cây Genre/Award khớp sơ đồ. Nhóm thể loại/giải được
ánh xạ theo nhãn trong bước build.

\textbf{Chuyển trang:} Sau đây nhóm chuyển sang r1 ·
contribu\allowbreak{}tion:\allowbreak{} who did what in a film?.

\subsection{\texorpdfstring{Slide 09 --- R1 ·
Contribu\allowbreak{}tion:\allowbreak{} who did what in a
film?}{Slide 09 --- R1 ·  who did what in a film?}}\label{slide-09-r1-contribution-who-did-what-in-a-film}

\textbf{Lời nói:}

Nếu chỉ gắn một nhãn nghề nghiệp cho Nolan, chúng ta không biết vai trò
đó thuộc phim nào. Contribution giải quyết bằng một bản ghi nối một
người, một phim và một vai trò. contribu\allowbreak{}tion\allowbreak{}By
trỏ Person, contribu\allowbreak{}tion\allowbreak{}To trỏ Film, hasRole
trỏ
Contribu\allowbreak{}tion\allowbreak{}Role\allowbreak{}.\allowbreak{}
Mỗi quan hệ có ràng buộc đúng một giá trị phù hợp. Bốn cá thể vai trò là
Director\allowbreak{}\allowbreak{}Role, ActorRole, WriterRole,
Producer\allowbreak{}\allowbreak{}Role.\allowbreak{} Nolan giữ ba vai
trò trong Inception, vì vậy có ba Contribution riêng. Các nhóm Acting,
Directing, Writing,
Producin\allowbreak{}g\allowbreak{}Contrib\allowbreak{}ution được phân
loại từ hasRole. Ảnh cần chụp Class Description của
Contribu\allowbreak{}tion, hiện cả exactly 1 và các restriction về loại
giá trị.

\textbf{Ảnh cần bổ sung --- P04:}
\texttt{P04\_\allowbreak{}cont\allowbreak{}ribution\allowbreak{}\_\allowbreak{}restric\allowbreak{}tions.\allowbreak{}pn\allowbreak{}g};
mở
\texttt{ontology\allowbreak{}/\allowbreak{}\allowbreak{}Movie\_\allowbreak{}\allowbreak{}O\allowbreak{}ntology.\allowbreak{}owl}.

\begin{enumerate}
\def\labelenumi{\arabic{enumi}.}
\tightlist
\item
  Entities → Classes → chọn Contribu\allowbreak{}tion.\allowbreak{}
\item
  Trong Class Description, mở đủ SubClass Of.
\item
  Hiện contribu\allowbreak{}tion\allowbreak{}By /
  contribu\allowbreak{}tion\allowbreak{}To / hasRole: exactly 1.
\end{enumerate}

\textbf{Cần thấy:} Ba qualified cardinality đúng 1, cùng các
all\allowbreak{}Value\allowbreak{}s\allowbreak{}From.\allowbreak{} Đây
là mô hình OWL, không phải biên bản kiểm tra thiếu trường.

\textbf{Chuyển trang:} Sau đây nhóm chuyển sang r1 · properties,
domain/range and inverses.

\subsection{Slide 10 --- R1 · Properties, domain/range and
inverses}\label{slide-10-r1-properties-domainrange-and-inverses}

\textbf{Lời nói:}

Object property nối hai thực thể.
contribu\allowbreak{}tion\allowbreak{}By có domain Contribution và range
Person; quan hệ ngược là
has\allowbreak{}Contr\allowbreak{}ibution.\allowbreak{}
contribu\allowbreak{}tion\allowbreak{}To nối tới Film, hasRole nối tới
Contribu\allowbreak{}tion\allowbreak{}Role\allowbreak{}.\allowbreak{} Ba
quan hệ này có tính functional. Các quan hệ trực tiếp director,
starring, writer, producer dùng từ DBpedia để hỏi thuận tiện. Những cặp
như
has\allowbreak{}Genre\allowbreak{}/\allowbreak{}genre\allowbreak{}Of\allowbreak{},
has\allowbreak{}Produ\allowbreak{}ction\allowbreak{}Com\allowbreak{}pany/\allowbreak{}pro\allowbreak{}duction\allowbreak{}O\allowbreak{}f,
has\allowbreak{}Award\allowbreak{}/\allowbreak{}award\allowbreak{}Of\allowbreak{}
cho phép đi hai chiều. inverse là đảo hướng, không đồng nghĩa quan hệ
đối xứng. Domain và range có thể giúp suy ra kiểu khi quan hệ được biết;
chúng không tự đóng vai trò từ chối nhập liệu như một biểu mẫu. Ảnh yêu
cầu chọn contribu\allowbreak{}tion\allowbreak{}By để thấy
domain/\allowbreak{}r\allowbreak{}ange/\allowbreak{}inv\allowbreak{}erse/\allowbreak{}fun\allowbreak{}ctional.\allowbreak{}

\textbf{Ảnh cần bổ sung --- P05:}
\texttt{P05\_\allowbreak{}obje\allowbreak{}ct\_\allowbreak{}prope\allowbreak{}rty.\allowbreak{}png};
mở
\texttt{ontology\allowbreak{}/\allowbreak{}\allowbreak{}Movie\_\allowbreak{}\allowbreak{}O\allowbreak{}ntology.\allowbreak{}owl}.

\begin{enumerate}
\def\labelenumi{\arabic{enumi}.}
\tightlist
\item
  Entities → Object properties → chọn
  contribu\allowbreak{}tion\allowbreak{}By.\allowbreak{}
\item
  Hiện Domain=Contribution, Range=Person.
\item
  Giữ inverse has\allowbreak{}Contr\allowbreak{}ibution và ô Functional
  được chọn.
\end{enumerate}

\textbf{Cần thấy:} Đúng
domain/\allowbreak{}r\allowbreak{}ange/\allowbreak{}inv\allowbreak{}erse/\allowbreak{}fun\allowbreak{}ctional
của contribu\allowbreak{}tion\allowbreak{}By.\allowbreak{} Nếu cần, chụp
bổ sung hasRole hoặc director.

\textbf{Tra cứu khi bảo vệ (không đọc toàn bộ):}

Reference: all 23 object properties (not intended to be read aloud in
full): \textbar{} Property \textbar{} Domain \textbar{} Range \textbar{}
Inverse / characte\allowbreak{}ristics \textbar{}
\textbar:--\textbar:--\textbar:--\textbar:--\textbar{} \textbar{}
dbo:director \textbar{} dbo:Film \textbar{} dbo:Person \textbar{}
ex:directed \textbar{} \textbar{} dbo:producer \textbar{} dbo:Film
\textbar{} dbo:Person \textbar{} ex:produced \textbar{} \textbar{}
dbo:starring \textbar{} dbo:Film \textbar{} dbo:Person \textbar{}
ex:actedIn \textbar{} \textbar{} dbo:writer \textbar{} dbo:Film
\textbar{} dbo:Person \textbar{} ex:wrote \textbar{} \textbar{}
ex:actedIn \textbar{} dbo:Person \textbar{} dbo:Film \textbar{}
dbo:starring \textbar{} \textbar{} ex:awardOf \textbar{} ex:Award
\textbar{} dbo:Film or dbo:Person \textbar{} ex:hasAward \textbar{}
\textbar{}
ex:\allowbreak{}contr\allowbreak{}ibution\allowbreak{}B\allowbreak{}y
\textbar{} ex:\allowbreak{}\allowbreak{}Contr\allowbreak{}ibution
\textbar{} dbo:Person \textbar{}
ex:\allowbreak{}has\allowbreak{}Co\allowbreak{}ntributi\allowbreak{}on ·
functional \textbar{} \textbar{}
ex:\allowbreak{}contr\allowbreak{}ibution\allowbreak{}O\allowbreak{}f
\textbar{} dbo:Film \textbar{}
ex:\allowbreak{}\allowbreak{}Contr\allowbreak{}ibution \textbar{}
ex:\allowbreak{}contr\allowbreak{}ibution\allowbreak{}T\allowbreak{}o
\textbar{} \textbar{}
ex:\allowbreak{}contr\allowbreak{}ibution\allowbreak{}T\allowbreak{}o
\textbar{} ex:\allowbreak{}\allowbreak{}Contr\allowbreak{}ibution
\textbar{} dbo:Film \textbar{}
ex:\allowbreak{}contr\allowbreak{}ibution\allowbreak{}O\allowbreak{}f ·
functional \textbar{} \textbar{} ex:country \textbar{} dbo:Film
\textbar{} dbo:Country \textbar{} --- \textbar{} \textbar{} ex:directed
\textbar{} dbo:Person \textbar{} dbo:Film \textbar{} dbo:director
\textbar{} \textbar{} ex:genreOf \textbar{} ex:Genre \textbar{} dbo:Film
\textbar{} ex:hasGenre \textbar{} \textbar{} ex:hasAward \textbar{}
dbo:Film or dbo:Person \textbar{} ex:Award \textbar{} ex:awardOf
\textbar{} \textbar{}
ex:\allowbreak{}has\allowbreak{}Co\allowbreak{}ntributi\allowbreak{}on
\textbar{} dbo:Person \textbar{}
ex:\allowbreak{}\allowbreak{}Contr\allowbreak{}ibution \textbar{}
ex:\allowbreak{}contr\allowbreak{}ibution\allowbreak{}B\allowbreak{}y
\textbar{} \textbar{} ex:hasGenre \textbar{} dbo:Film \textbar{}
ex:Genre \textbar{} ex:genreOf \textbar{} \textbar{}
ex:\allowbreak{}has\allowbreak{}Pr\allowbreak{}oduction\allowbreak{}\allowbreak{}Company
\textbar{} dbo:Film \textbar{}
ex:\allowbreak{}\allowbreak{}Produ\allowbreak{}ction\allowbreak{}Com\allowbreak{}pany
\textbar{} ex:\allowbreak{}produ\allowbreak{}ction\allowbreak{}Of
\textbar{} \textbar{} ex:hasRole \textbar{}
ex:\allowbreak{}\allowbreak{}Contr\allowbreak{}ibution \textbar{}
ex:\allowbreak{}\allowbreak{}Contr\allowbreak{}ibution\allowbreak{}R\allowbreak{}ole
\textbar{} ex:roleOf · functional \textbar{} \textbar{} ex:language
\textbar{} dbo:Film \textbar{} ex:Language \textbar{} --- \textbar{}
\textbar{} ex:produced \textbar{} dbo:Person \textbar{} dbo:Film
\textbar{} dbo:producer \textbar{} \textbar{}
ex:\allowbreak{}produ\allowbreak{}ction\allowbreak{}Of \textbar{}
ex:\allowbreak{}\allowbreak{}Produ\allowbreak{}ction\allowbreak{}Com\allowbreak{}pany
\textbar{} dbo:Film \textbar{}
ex:\allowbreak{}has\allowbreak{}Pr\allowbreak{}oduction\allowbreak{}\allowbreak{}Company
\textbar{} \textbar{} ex:roleOf \textbar{}
ex:\allowbreak{}\allowbreak{}Contr\allowbreak{}ibution\allowbreak{}R\allowbreak{}ole
\textbar{} ex:\allowbreak{}\allowbreak{}Contr\allowbreak{}ibution
\textbar{} ex:hasRole \textbar{} \textbar{}
ex:\allowbreak{}sourc\allowbreak{}e\allowbreak{}Snapsho\allowbreak{}t
\textbar{} --- \textbar{}
ex:\allowbreak{}\allowbreak{}Sourc\allowbreak{}e\allowbreak{}Snapsho\allowbreak{}t
\textbar{} --- \textbar{} \textbar{} ex:wrote \textbar{} dbo:Person
\textbar{} dbo:Film \textbar{} dbo:writer \textbar{}

\textbf{Chuyển trang:} Sau đây nhóm chuyển sang r1 · how owl expresses
meaning.

\subsection{Slide 11 --- R1 · How OWL expresses
meaning}\label{slide-11-r1-how-owl-expresses-meaning}

\textbf{Lời nói:}

OWL cho phép định nghĩa tương đương lớp, giao, hợp, tồn tại, giá trị cố
định và cardinality. Ví dụ Contribution có exactly một Person, một Film
và một Role;
Directin\allowbreak{}g\allowbreak{}Contrib\allowbreak{}ution được nhận
diện khi hasRole có giá trị
Director\allowbreak{}\allowbreak{}Role.\allowbreak{} Quan hệ inverse suy
ra đường đi ngược. Disjointness chỉ dùng cho các lớp không thể chồng lấp
theo mô hình, không dùng cho Actor với Filmmaker. Cần nhớ giả định thế
giới mở: thiếu một triple không tự chứng minh sự việc không có. Các IRI
khác nhau cũng không tự bảo đảm là cá thể khác nhau. Vì vậy ứng dụng
dùng kiểm tra Python riêng để phát hiện thiếu hoặc thừa trường, còn OWL
phát biểu ý nghĩa và hỗ trợ phân loại.

\textbf{Chuyển trang:} Sau đây nhóm chuyển sang r1 · fourteen classes
with additional types.

\subsection{Slide 12 --- R1 · Fourteen classes with additional
types}\label{slide-12-r1-fourteen-classes-with-additional-types}

\textbf{Lời nói:}

Bảng này cho thấy toàn bộ 14 lớp có kiểu được bổ sung. Bốn nhóm
Contribution dùng role, ba nhóm người dùng đóng góp hoặc giải, năm nhóm
phim dùng thể loại hoặc giải. Mười hai lớp thuộc phần xử lý bằng luật
OWL RL. Hai lớp còn lại,
Multi\allowbreak{}Gen\allowbreak{}re\allowbreak{}Film và FilmStudio,
dùng truy vấn COUNT DISTINCT sau bước luật. Các số đếm chỉ tính tài
nguyên nội bộ để tránh alias sameAs làm tăng số. Bảng là kết quả trong
ontology\allowbreak{}\_\allowbreak{}reasoni\allowbreak{}ng.\allowbreak{}json,\allowbreak{}
không phải 14 kiểu gán sẵn trong movies.ttl. Hai truy vấn người có LIMIT
20; số hiển thị không phải tổng 769 Actor hoặc 89 Filmmaker. Phần
cardinality được giải thích riêng ở slide sau.

\textbf{Chuyển trang:} Sau đây nhóm chuyển sang r1 · inference example:
christopher nolan.

\subsection{Slide 13 --- R1 · Inference example: Christopher
Nolan}\label{slide-13-r1-inference-example-christopher-nolan}

\textbf{Lời nói:}

Dữ liệu gốc chỉ khai báo Nolan là Person, cùng các quan hệ đóng góp và
giải. Một Contribution của Nolan trong Inception có
Director\allowbreak{}\allowbreak{}Role.\allowbreak{} Định nghĩa hasValue
phân loại bản ghi này là
Directin\allowbreak{}g\allowbreak{}Contrib\allowbreak{}ution.\allowbreak{}
Từ Person có loại đóng góp đạo diễn, biên kịch hoặc sản xuất, định nghĩa
dùng giao, hợp và tồn tại phân loại Nolan là Filmmaker. Nolan còn có
giải riêng nên được bổ sung AwardWinner. Đây là khác biệt giữa kiểu gán
ban đầu và kiểu suy ra trong ứng dụng. Khung ảnh yêu cầu Class
Description và Equivalent To của Filmmaker. Ảnh này chứng minh công thức
đã khai báo; kết quả thực thi được kiểm chứng riêng bằng terminal và
tests.

\textbf{Ảnh cần bổ sung --- P06:}
\texttt{P06\_\allowbreak{}film\allowbreak{}maker\_\allowbreak{}eq\allowbreak{}uivalent\allowbreak{}\_\allowbreak{}class.\allowbreak{}p\allowbreak{}ng};
mở
\texttt{ontology\allowbreak{}/\allowbreak{}\allowbreak{}Movie\_\allowbreak{}\allowbreak{}O\allowbreak{}ntology.\allowbreak{}owl}.

\begin{enumerate}
\def\labelenumi{\arabic{enumi}.}
\tightlist
\item
  Entities → Classes → chọn Filmmaker.
\item
  Mở Class Description → Equivalent To.
\item
  Hiện Person AND các nhánh SOME
  Directin\allowbreak{}g/\allowbreak{}\allowbreak{}Writin\allowbreak{}g/\allowbreak{}\allowbreak{}Produc\allowbreak{}ing\allowbreak{}Contr\allowbreak{}ibution.\allowbreak{}
\end{enumerate}

\textbf{Cần thấy:} Định nghĩa giao/hợp/tồn tại. Ảnh định nghĩa không tự
chứng minh HermiT đã phân loại dữ liệu.

\textbf{Chuyển trang:} Sau đây nhóm chuyển sang r1 · cardinality:
thresholds of two and three.

\subsection{Slide 14 --- R1 · Cardinality: thresholds of two and
three}\label{slide-14-r1-cardinality-thresholds-of-two-and-three}

\textbf{Lời nói:}

Multi\allowbreak{}Gen\allowbreak{}re\allowbreak{}Film được định nghĩa là
Film có ít nhất hai Genre; FilmStudio là
Producti\allowbreak{}on\allowbreak{}Compan\allowbreak{}y có ít nhất ba
Film qua producti\allowbreak{}on\allowbreak{}Of.\allowbreak{} Trong mẫu,
toàn bộ 30 phim có ít nhất hai thể loại và có sáu công ty đạt ngưỡng ba
phim. Warner Bros.~liên quan tới 11 phim trong mẫu. owlrl không xử lý
các ngưỡng này theo nhu cầu của bài, nên reason.py dùng COUNT DISTINCT
bổ sung sau khi suy ra inverse. Đếm IRI khác nhau là quy tắc ứng dụng;
trong OWL, hai tên khác nhau có thể chỉ cùng một cá thể.
Hermi\allowbreak{}T/\allowbreak{}\allowbreak{}P\allowbreak{}ellet đã
chạy: hai lớp này có không cá thể được suy luận, còn 30 và sáu là số đếm
trong ứng dụng. Ảnh cần chụp công thức min 3 của FilmStudio, không yêu
cầu giả lập kết quả reasoner.

\textbf{Ảnh cần bổ sung --- P07:}
\texttt{P07\_\allowbreak{}film\allowbreak{}studio\_\allowbreak{}c\allowbreak{}ardinali\allowbreak{}ty.\allowbreak{}png};
mở
\texttt{ontology\allowbreak{}/\allowbreak{}\allowbreak{}Movie\_\allowbreak{}\allowbreak{}O\allowbreak{}ntology.\allowbreak{}owl}.

\begin{enumerate}
\def\labelenumi{\arabic{enumi}.}
\tightlist
\item
  Entities → Classes → chọn FilmStudio.
\item
  Mở Equivalent To; hiện min 3 productionOf Film.
\item
  Có thể chụp thêm
  Multi\allowbreak{}Gen\allowbreak{}re\allowbreak{}Film:\allowbreak{}
  min 2 hasGenre Genre.
\end{enumerate}

\textbf{Cần thấy:} Ảnh công thức cardinality. Kết quả 6 studio/30 phim
nhiều thể loại của app dùng COUNT DISTINCT, không tự coi là chứng minh
OWL DL.

\textbf{Chuyển trang:} Sau đây nhóm chuyển sang r2 · real collection and
traceable sources.

\subsection{Slide 15 --- R2 · Real collection and traceable
sources}\label{slide-15-r2-real-collection-and-traceable-sources}

\textbf{Lời nói:}

Danh tính phim được xác định qua sitelink Wikipedia tiếng Anh chính xác,
thay vì so tên gần giống. Wikidata cung cấp đạo diễn P57, diễn viên
P161, biên kịch P58, nhà sản xuất P162, giải P166 và công ty P272.
DBpedia chỉ nối khi chủ thể khớp tiêu đề và có kiểu Film; vì vậy có 28
liên kết DBpedia thay vì tự đoán đủ 30. Mỗi phản hồi có URL, provider,
thời điểm, HTTP status, SHA-256 và đường dẫn. Hiện đủ 76 phản hồi và
hash khớp. Một phản hồi tải lại đổi nội dung được ghi hash/time mới,
đồng thời giữ danh mục lịch sử. Hash kiểm tra toàn vẹn byte, không chứng
minh mọi phát biểu ngoài đời đúng.

\textbf{Chuyển trang:} Sau đây nhóm chuyển sang r3 · rdf and typed
literal values.

\subsection{Slide 16 --- R3 · RDF and typed literal
values}\label{slide-16-r3-rdf-and-typed-literal-values}

\textbf{Lời nói:}

Một triple gồm chủ thể, quan hệ và đối tượng hoặc giá trị. Inception có
đạo diễn Nolan, năm 2010 và thời lượng 148 phút. Quan hệ đạo diễn nối
hai thực thể; năm và thời lượng là literal có datatype. Dùng QID để tạo
IRI ổn định; chuẩn hóa các đơn vị giờ, phút, giây về phút và ghi các
trường hợp nguồn có nhiều giá trị. Sáu datatype property khai báo là
title, releaseYear, runtime\allowbreak{}M\allowbreak{}inutes, sourceUrl,
retrievedAt, sha256. Turtle và JSON-LD biểu diễn cùng graph; file OWL
riêng có schema, file Knowledge Graph có schema và dữ liệu. Ảnh cần chọn
runtime\allowbreak{}M\allowbreak{}inutes trong Data properties để thấy
domain Film và range decimal.

\textbf{Ảnh cần bổ sung --- P08:}
\texttt{P08\_\allowbreak{}runt\allowbreak{}ime\_\allowbreak{}data\allowbreak{}type.\allowbreak{}png\allowbreak{}};
mở
\texttt{ontology\allowbreak{}/\allowbreak{}\allowbreak{}Movie\_\allowbreak{}\allowbreak{}O\allowbreak{}ntology.\allowbreak{}owl}.

\begin{enumerate}
\def\labelenumi{\arabic{enumi}.}
\tightlist
\item
  Entities → Data properties → chọn
  runtime\allowbreak{}M\allowbreak{}inutes.\allowbreak{}
\item
  Hiện Domain=Film và Range=xsd:decimal.
\item
  Có thể chụp thêm releaseYear có range xsd:integer.
\end{enumerate}

\textbf{Cần thấy:} Thuộc tính dữ liệu nối thực thể với literal, khác
quan hệ nối hai thực thể.

\textbf{Chuyển trang:} Sau đây nhóm chuyển sang r3--r4 · from rdf to
linked open data.

\subsection{Slide 17 --- R3--R4 · From RDF to Linked Open
Data}\label{slide-17-r3r4-from-rdf-to-linked-open-data}

\textbf{Lời nói:}

Mức một sao bắt đầu với dữ liệu trên Web có giấy phép mở; hai sao là có
cấu trúc; ba sao dùng định dạng mở; bốn sao dùng HTTP URI và chuẩn RDF
để tra cứu; năm sao thêm liên kết tới dữ liệu khác. Bài có giấy phép CC
BY-SA 4.0,
R\allowbreak{}D\allowbreak{}F/\allowbreak{}\allowbreak{}Turt\allowbreak{}le/\allowbreak{}\allowbreak{}J\allowbreak{}S\allowbreak{}O\allowbreak{}N-\allowbreak{}\allowbreak{}L\allowbreak{}D
và các IRI có trang mô tả. Bản công khai đã được kiểm tra không đăng
nhập, cả bốn graph RDF đối chiếu đều đẳng cấu với local. Có 1.699 liên
kết Wikidata và 28 DBpedia. sameAs khẳng định cùng danh tính, còn
source\allowbreak{}Sn\allowbreak{}apshot ghi xuất xứ. Tái dùng dbo:Film
là dùng từ vựng lớp ở YC1, khác liên kết cá thể ở YC4.

\textbf{Chuyển trang:} Sau đây nhóm chuyển sang r3 · looking up
inception's iri.

\subsection{Slide 18 --- R3 · Looking up Inception's
IRI}\label{slide-18-r3-looking-up-inceptions-iri}

\textbf{Lời nói:}

Trang Inception hiển thị tên, kiểu, năm, thời lượng, đạo diễn, đóng góp,
liên kết ngoài và nguồn. Người đọc có thể tải mô tả Turtle; HTML cũng
nhúng JSON-LD. Ở Flask local, Accept text turtle trả 303 đến mô tả RDF
và curl trừ L theo chuyển hướng. Hosted tĩnh cung cấp HTML/JSON-LD và
link RDF; không tự khẳng định nó có cùng hành vi 303 như Flask. Bên trái
là ảnh ứng dụng thật. Bên phải là khung cần ảnh Individuals trong
Protégé khi mở Knowledge Graph: chọn film-Q25188 và hiện kiểu cùng các
thuộc tính. Không mở file chỉ schema để chụp cá thể.

\textbf{Ảnh cần bổ sung --- P09:}
\texttt{P09\_\allowbreak{}ince\allowbreak{}ption\_\allowbreak{}in\allowbreak{}dividual\allowbreak{}.\allowbreak{}png};
mở
\texttt{ontology\allowbreak{}/\allowbreak{}\allowbreak{}Movie\_\allowbreak{}\allowbreak{}K\allowbreak{}nowledge\allowbreak{}\_\allowbreak{}\allowbreak{}Graph.\allowbreak{}o\allowbreak{}wl}.

\begin{enumerate}
\def\labelenumi{\arabic{enumi}.}
\tightlist
\item
  Mở
  ontology\allowbreak{}/\allowbreak{}\allowbreak{}Movie\_\allowbreak{}\allowbreak{}K\allowbreak{}nowledge\allowbreak{}\_\allowbreak{}\allowbreak{}Graph.\allowbreak{}o\allowbreak{}wl.\allowbreak{}
\item
  Individuals → chọn film-Q25188 (Inception).
\item
  Hiện Film, title, releaseYear=2010, runtimeMinutes=148 và director.
\end{enumerate}

\textbf{Cần thấy:} Cá thể thật và thuộc tính khớp truy vấn. Giữ tên
ontology để không chụp nhầm file chỉ có schema.

\textbf{Chuyển trang:} Sau đây nhóm chuyển sang r5 · querying inception
on the web.

\subsection{Slide 19 --- R5 · Querying Inception on the
Web}\label{slide-19-r5-querying-inception-on-the-web}

\textbf{Lời nói:}

Truy vấn đầu tiên tìm phim tên Inception, đi theo director tới người và
lấy nhãn người đó. Kết quả là Inception, 2010, 148 phút, Christopher
Nolan. SELECT chọn cột, WHERE đưa mẫu đồ thị, dấu hỏi đánh dấu biến.
OPTIONAL giữ dòng khi thuộc tính năm hoặc thời lượng có thể thiếu. Người
dùng có thể sửa truy vấn, bấm Run query và tải kết quả. Giao diện có 14
mẫu trực tiếp. Ngoài SELECT trả bảng, ASK trả True hoặc False, CONSTRUCT
tạo graph RDF. Kết quả được tính từ graph; danh sách phim hoặc tên đạo
diễn không được nhập cứng trong bảng.

\textbf{Chuyển trang:} Sau đây nhóm chuyển sang r5 · one person,
multiple roles and questions.

\subsection{Slide 20 --- R5 · One person, multiple roles and
questions}\label{slide-20-r5-one-person-multiple-roles-and-questions}

\textbf{Lời nói:}

Câu hỏi toàn bộ đóng góp Inception trả 25 bản ghi, gồm 21 diễn viên, một
đạo diễn, một biên kịch và hai nhà sản xuất. Ảnh trên slide lọc riêng
Christopher Nolan nên chỉ còn ba dòng: Director, Writer, Producer. Chúng
ta không lấy ba dòng này làm tổng đóng góp. Các truy vấn khác cho biết
tám phim do Nolan đạo diễn trong mẫu, bốn công ty liên quan Inception và
bảy giải của The Godfather. Việc thêm câu hỏi không cần nhập một bảng
kết quả riêng; chúng dùng cùng đồ thị và mô hình. Download results xuất
kết quả đang hiển thị. Các file truy vấn và biên bản cho phép kiểm tra
lại những số này.

\textbf{Chuyển trang:} Sau đây nhóm chuyển sang r5 · endpoint, terminal
and inferred types.

\subsection{Slide 21 --- R5 · Endpoint, terminal and inferred
types}\label{slide-21-r5-endpoint-terminal-and-inferred-types}

\textbf{Lời nói:}

Bài có endpoint GET/POST và terminal, ngoài giao diện Web. SELECT và ASK
trả JSON; CONSTRUCT và DESCRIBE trả Turtle. Endpoint mặc định chỉ nạp
graph khai báo. Lệnh query.py với reasoned nạp thêm schema và
inferred\allowbreak{}\_\allowbreak{}classes\allowbreak{}.\allowbreak{}ttl.\allowbreak{}
Cùng câu 24 trên Nolan trả một kiểu Person khi chạy gốc; chạy reasoned
trả Person, Filmmaker, AwardWinner. File kiểu được tạo trước bởi
reason.py, không tính suy luận mỗi lần gọi query.py. Truy vấn hierarchy
như câu 13 cần schema;
Actor/\allowbreak{}\allowbreak{}Fi\allowbreak{}lmmaker có LIMIT 20 nên
tổng thực ở biên bản là 769 và 89. Khung ảnh Protégé yêu cầu Types của
Nolan ở chế độ asserted; nếu ảnh sau reasoner phải ghi rõ trạng thái.

\textbf{Ảnh cần bổ sung --- P10:}
\texttt{P10\_\allowbreak{}nola\allowbreak{}n\_\allowbreak{}assert\allowbreak{}ed\_\allowbreak{}types\allowbreak{}.\allowbreak{}png};
mở
\texttt{ontology\allowbreak{}/\allowbreak{}\allowbreak{}Movie\_\allowbreak{}\allowbreak{}K\allowbreak{}nowledge\allowbreak{}\_\allowbreak{}\allowbreak{}Graph.\allowbreak{}o\allowbreak{}wl}.

\begin{enumerate}
\def\labelenumi{\arabic{enumi}.}
\tightlist
\item
  Mở
  Movie\_\allowbreak{}\allowbreak{}Kn\allowbreak{}owledge\_\allowbreak{}\allowbreak{}Graph.\allowbreak{}ow\allowbreak{}l,
  chọn person-\allowbreak{}\allowbreak{}Q\allowbreak{}25191 (Nolan).
\item
  Hiện Types trong chế độ asserted / trước suy luận.
\item
  Giữ dbo:Person và các quan hệ đóng góp/giải trong ảnh.
\end{enumerate}

\textbf{Cần thấy:} Kiểu gán gốc là Person;
Filmmake\allowbreak{}r/\allowbreak{}\allowbreak{}Award\allowbreak{}W\allowbreak{}inner
trong app được bổ sung qua file phân loại. Nếu chụp sau reasoner, phải
ghi rõ tên và trạng thái thực thi.

\textbf{Tra cứu khi bảo vệ (không đọc toàn bộ):}

Reference: 24 queries; values are result-row counts or ASK booleans:
\textbar{} File \textbar{} Asserted \textbar{} With schema / inference
\textbar{} \textbar:--\textbar:--\textbar:--\textbar{} \textbar{}
01\_films.rq \textbar{} 30 \textbar{} 30 \textbar{} \textbar{}
02\_\allowbreak{}incep\allowbreak{}tion.\allowbreak{}rq \textbar{} 1
\textbar{} 1 \textbar{} \textbar{} 03\_nolan.rq \textbar{} 8 \textbar{}
8 \textbar{} \textbar{}
04\_\allowbreak{}credi\allowbreak{}ts.\allowbreak{}rq \textbar{} 25
\textbar{} 25 \textbar{} \textbar{}
05\_\allowbreak{}exter\allowbreak{}nal\_\allowbreak{}link\allowbreak{}s.\allowbreak{}rq
\textbar{} 58 \textbar{} 58 \textbar{} \textbar{} 06\_genres.rq
\textbar{} 75 \textbar{} 75 \textbar{} \textbar{} 07\_source.rq
\textbar{} 2 \textbar{} 2 \textbar{} \textbar{} 08\_ask.rq \textbar{}
true \textbar{} true \textbar{} \textbar{}
09\_\allowbreak{}actor\allowbreak{}s\_\allowbreak{}of\_\allowbreak{}inc\allowbreak{}eption.\allowbreak{}r\allowbreak{}q
\textbar{} 21 \textbar{} 21 \textbar{} \textbar{}
10\_\allowbreak{}produ\allowbreak{}ction\_\allowbreak{}co\allowbreak{}mpanies\_\allowbreak{}of\_\allowbreak{}a\_\allowbreak{}fil\allowbreak{}m.\allowbreak{}rq
\textbar{} 4 \textbar{} 4 \textbar{} \textbar{}
11\_\allowbreak{}award\allowbreak{}s\_\allowbreak{}receiv\allowbreak{}ed\_\allowbreak{}by\_\allowbreak{}a\_\allowbreak{}film.\allowbreak{}rq
\textbar{} 7 \textbar{} 7 \textbar{} \textbar{}
12\_\allowbreak{}count\allowbreak{}ries\_\allowbreak{}and\allowbreak{}\_\allowbreak{}languag\allowbreak{}es\_\allowbreak{}of\_\allowbreak{}a\_\allowbreak{}film.\allowbreak{}rq
\textbar{} 1 \textbar{} 1 \textbar{} \textbar{}
13\_\allowbreak{}ficti\allowbreak{}on\_\allowbreak{}genre\allowbreak{}\_\allowbreak{}films\_\allowbreak{}v\allowbreak{}ia\_\allowbreak{}hiera\allowbreak{}rchy.\allowbreak{}rq
\textbar{} 0 \textbar{} 29 \textbar{} \textbar{}
14\_\allowbreak{}all\_\allowbreak{}a\allowbreak{}gents\_\allowbreak{}pe\allowbreak{}ople\_\allowbreak{}and\allowbreak{}\_\allowbreak{}organiz\allowbreak{}ations.\allowbreak{}r\allowbreak{}q
\textbar{} 1 \textbar{} 1 \textbar{} \textbar{}
15\_\allowbreak{}featu\allowbreak{}re\_\allowbreak{}vs\_\allowbreak{}an\allowbreak{}imated\_\allowbreak{}f\allowbreak{}ilm\_\allowbreak{}coun\allowbreak{}ts.\allowbreak{}rq
\textbar{} 1 \textbar{} 1 \textbar{} \textbar{}
16\_\allowbreak{}docum\allowbreak{}entary\_\allowbreak{}f\allowbreak{}ilm\_\allowbreak{}coun\allowbreak{}t\_\allowbreak{}ask.\allowbreak{}rq\allowbreak{}
\textbar{} false \textbar{} false \textbar{} \textbar{}
17\_\allowbreak{}infer\allowbreak{}red\_\allowbreak{}acto\allowbreak{}rs.\allowbreak{}rq
\textbar{} 0 \textbar{} 20 \textbar{} \textbar{}
18\_\allowbreak{}infer\allowbreak{}red\_\allowbreak{}film\allowbreak{}makers.\allowbreak{}r\allowbreak{}q
\textbar{} 0 \textbar{} 20 \textbar{} \textbar{}
19\_\allowbreak{}infer\allowbreak{}red\_\allowbreak{}acti\allowbreak{}on\_\allowbreak{}films\allowbreak{}.\allowbreak{}rq
\textbar{} 0 \textbar{} 12 \textbar{} \textbar{}
20\_\allowbreak{}infer\allowbreak{}red\_\allowbreak{}mult\allowbreak{}i\_\allowbreak{}genre\_\allowbreak{}films.\allowbreak{}rq\allowbreak{}
\textbar{} 0 \textbar{} 30 \textbar{} \textbar{}
21\_\allowbreak{}infer\allowbreak{}red\_\allowbreak{}awar\allowbreak{}d\_\allowbreak{}winnin\allowbreak{}g\_\allowbreak{}films.\allowbreak{}rq
\textbar{} 0 \textbar{} 26 \textbar{} \textbar{}
22\_\allowbreak{}infer\allowbreak{}red\_\allowbreak{}film\allowbreak{}\_\allowbreak{}studios\allowbreak{}.\allowbreak{}rq
\textbar{} 0 \textbar{} 6 \textbar{} \textbar{}
23\_\allowbreak{}peopl\allowbreak{}e\_\allowbreak{}with\_\allowbreak{}d\allowbreak{}irecting\allowbreak{}\_\allowbreak{}and\_\allowbreak{}wri\allowbreak{}ting\_\allowbreak{}con\allowbreak{}tributio\allowbreak{}n.\allowbreak{}rq
\textbar{} 0 \textbar{} 10 \textbar{} \textbar{}
24\_\allowbreak{}nolan\allowbreak{}\_\allowbreak{}asserte\allowbreak{}d\_\allowbreak{}types\_\allowbreak{}only.\allowbreak{}rq
\textbar{} 1 \textbar{} 3 \textbar{} Queries 14/15 each return one
aggregate row. Queries 17/18 have LIMIT 20; full
Actor/\allowbreak{}\allowbreak{}Fi\allowbreak{}lmmaker totals are
769/89. Query 24 returns one asserted type or three with inference.

\textbf{Chuyển trang:} Sau đây nhóm chuyển sang validating data,
application and publication.

\subsection{Slide 22 --- Validating data, application and
publication}\label{slide-22-validating-data-application-and-publication}

\textbf{Lời nói:}

validate đã kiểm tra tên, nguồn, liên kết của phim và các thành phần
Contribu\allowbreak{}tion; 76 hash nguồn khớp và 24 file truy vấn đã
chạy ở chế độ dữ liệu phù hợp. Có 14 test pass và chín kiểm tra trình
duyệt đạt, gồm endpoint, Comunica, lỗi cú pháp, ASK, CONSTRUCT, trang
IRI và màn hình mobile. Bản công khai được kiểm tra qua tám URL không có
cookie hay đăng nhập; bốn graph RDF đều đẳng cấu với local, có giấy phép
và link mô tả máy đọc được.
Hermi\allowbreak{}T/\allowbreak{}\allowbreak{}P\allowbreak{}ellet riêng
đã xác nhận OWL nhất quán và không có lớp không khả thỏa. Test và hash
không chứng minh mọi thông tin ngoài đời đúng. Timestamp
O\allowbreak{}W\allowbreak{}L/\allowbreak{}\allowbreak{}Turt\allowbreak{}le/\allowbreak{}\allowbreak{}J\allowbreak{}S\allowbreak{}O\allowbreak{}N-\allowbreak{}\allowbreak{}L\allowbreak{}D
đã đồng bộ đến mili giây. Đó là lý do cần vừa nêu kết quả vừa giữ giới
hạn phương pháp.

\textbf{Chuyển trang:} Sau đây nhóm chuyển sang assessment and
limitations.

\subsection{Slide 23 --- Assessment and
limitations}\label{slide-23-assessment-and-limitations}

\textbf{Lời nói:}

Ảnh đề gốc không đưa trọng số; nhóm dùng thang tự đánh giá chia đều hai
điểm mỗi yêu cầu. Ontology, thu thập, RDF công khai, liên kết và truy
vấn đều có minh chứng đạt nên đề xuất mười trên mười. Đây không phải
điểm chính thức của giảng viên. Các giới hạn vẫn còn: mẫu 30 phim có chủ
đích; nhóm thể loại và giải ánh xạ theo nhãn; chưa có ngân sách hoặc
streaming;
Hermi\allowbreak{}T/\allowbreak{}\allowbreak{}P\allowbreak{}ellet đã xác
nhận nhất quán, còn hai nhóm cardinality chỉ có số đếm ứng dụng 30/6; DL
có không cá thể. SQL cũng có thể trả nhiều câu hỏi bằng
J\allowbreak{}O\allowbreak{}I\allowbreak{}N/\allowbreak{}vie\allowbreak{}w/\allowbreak{}quy
tắc. Giá trị Semantic Web ở đây là IRI dùng chung, từ vựng tái dùng,
liên kết dataset, xuất xứ và định nghĩa ngữ nghĩa tường minh.

\textbf{Chuyển trang:} Sau đây nhóm chuyển sang conclusion and pending
protégé evidence.

\subsection{Slide 24 --- Conclusion and pending Protégé
evidence}\label{slide-24-conclusion-and-pending-protuxe9guxe9-evidence}

\textbf{Lời nói:}

MovieLOD đã nối mô hình, dữ liệu thật, nguồn, RDF, liên kết và truy vấn
trong một ứng dụng. Bản này có 24 slide; nhóm có thể trình bày đầy đủ
khoảng 18--22 phút hoặc chọn các trang chính để rút gọn. Trước khi nộp,
điền tên thành viên và bổ sung các ảnh Protégé theo mã trên bảng. Các
khung ảnh và hướng dẫn đã có ngay ở những slide liên quan. Giao diện
macOS chưa cho phép chụp trong lần làm tài liệu này, nên không có ảnh
Protégé giả được gắn vào. Có thể chèn ảnh trực tiếp trong PowerPoint
hoặc lưu đúng tên ở
evidence\allowbreak{}/\allowbreak{}protege\allowbreak{} rồi chạy lại
trình tạo slide. Nhóm xin kết thúc và sẵn sàng trả lời câu hỏi.

\textbf{Chuyển trang:} Nhóm xin mời thầy cô đặt câu hỏi.

\subsection{Câu hỏi bảo vệ
ngắn}\label{cuxe2u-hux1ecfi-bux1ea3o-vux1ec7-ngux1eafn}

\textbf{42 lớp khác gì 30 phim?} 42 là lớp trong schema; 30 là cá thể
Film trong graph dữ liệu. Protégé Metrics có thể tính thêm lớp ngoài
được tham chiếu.

\textbf{Vì sao cần Contribution?} Một người có nhiều vai trò trong nhiều
phim; mỗi bộ người--phim--vai trò có bản ghi riêng.

\textbf{Vì sao exact cardinality chưa thay Python?} OWL dùng thế giới
mở; Python kiểm tra trường bắt buộc của ứng dụng.

\textbf{COUNT DISTINCT có phải suy luận DL không?} Đếm tên IRI là quy
tắc ứng dụng; OWL không mặc định tên khác nhau chỉ cá thể khác nhau.

\textbf{Endpoint không có Filmmaker có phải lỗi?} Endpoint mặc định
graph gốc. Dùng --reasoned nạp schema và kết quả phân loại.

\textbf{sameAs khác nguồn thế nào?} sameAs là cùng danh tính;
source\allowbreak{}Sn\allowbreak{}apshot là xuất xứ; dbo:Film là tái
dùng từ vựng.

\textbf{SQL có trả được các câu hỏi không?} Có, bằng
J\allowbreak{}O\allowbreak{}I\allowbreak{}N/\allowbreak{}vie\allowbreak{}w/\allowbreak{}quy
tắc. Giá trị của bài là IRI, từ vựng chung, liên kết, nguồn và định
nghĩa ngữ nghĩa.

\textbf{Ảnh Protégé có chứng minh reasoner đã chạy không?} Ảnh cây
asserted/định nghĩa chỉ chứng minh khai báo. Muốn dùng ảnh inferred,
phải ghi rõ reasoner và trạng thái/kết quả thật.

\end{document}
